% ===========================================================================
%  ASTROLABE blog post, editable source
%
%  Write the prose here; hand the file back and it gets pushed to
%  blog/astrolabe-flow-matching.html. Everything the page renders has a placeholder
%  below, in page order, so moving a \interactive or \figplace block moves
%  the real thing on the page.
%
%  Compile with pdflatex to see the layout with the placeholders drawn in.
%
%  Rules of the mapping:
%    \section{...}        -> an <h2> on the page
%    \subsection{...}     -> an <h3> on the page
%    ordinary paragraphs  -> <p>
%    \lead{...}           -> the opening paragraph, set larger
%    \interactive{key}{}  -> a live figure, keys listed below
%    \figplace{file}{}{}  -> a static figure: figures/<file>.png on the page,
%                            pdf/<file>.pdf behind its PDF link; [narrow]
%                            holds it to about half the column, [medium] a
%                            little under the whole
%    \videoplace{stem}{}  -> a video from haithem-assets/videos, played at
%                            the interactive figures' width
%    \figpair{a}{b}{}     -> two pictures side by side at the text's width,
%                            one caption under both
%    figrow environment   -> the figures inside share one row on the page
%    gallery environment  -> a grid of movie posters that open a player;
%                            \movie{group/stem}{title}{caption} inside it is
%                            one movie, from movies/ in the data repository
%    \pullquote{...}      -> a large centred quote between sections
%    \pullquote*{...}     -> the same panel without the quotation marks
%    \pullquote[b]{a}     -> two of them, "a" then an arrow then "b"
%    \wide                -> as the last argument, breaks the figure out of
%                            the text column (for very wide diagrams)
%    \astrolabe           -> a name, styled the same wherever it appears
%    \papercard            -> not text on the page: the paper's title, authors
%                            and abstract, which become the page's search
%                            metadata (description and JSON-LD)
%    thebibliography      -> the References list at the foot of the page,
%                            \cite keys become superscript links into it
%    \indexcard            -> the card this post gets on blog.html, the
%                            listing page. It is not part of the post; it
%                            sits at the end of this file, after the
%                            references, and is drawn last on paper too.
%
%  Interactive keys, do not rename:
%    astro3d    the 3D box: N-body truth vs a model through a
%               draggable split, with the movie's timeline on Play
% ===========================================================================

\documentclass[11pt]{article}
\usepackage[margin=1in]{geometry}
\usepackage{xcolor}
\usepackage{tcolorbox}
\usepackage{hyperref}
\usepackage{parskip}

% the name, in small caps wherever it appears
\newcommand{\astrolabe}{\textsc{astrolabe}}

\newcommand{\posttitle}[1]{\title{#1}\def\theposttitle{#1}}
\newcommand{\posttags}[1]{\def\theposttags{#1}}
\newcommand{\postdate}[1]{\date{#1}\def\thepostdate{#1}}
\newcommand{\postbyline}[1]{\def\thepostbyline{#1}}

% opening paragraph, rendered larger on the page
\newcommand{\lead}[1]{{\large #1}\par\medskip}

% a live figure on the page
\newcommand{\interactive}[3][]{%
  \begin{tcolorbox}[colback=orange!5,colframe=orange!60!black,
                    title={Interactive: \texttt{#2}\ifx\relax#1\relax\else\ (#1)\fi}]
  \emph{Caption:} #3
  \end{tcolorbox}}

% a static figure; #1 optional "wide", #2 file stem, #3 caption
\newcommand{\figplace}[3][]{%
  \begin{tcolorbox}[colback=blue!4,colframe=blue!40!black,
                    title={Figure: \texttt{#2}\ifx\relax#1\relax\else\ (#1)\fi}]
  \emph{Caption:} #3
  \end{tcolorbox}}

% a recorded figure: #1 file stem under haithem-assets/videos, #2 caption.
% It plays on the page at the interactive figures' width
\newcommand{\videoplace}[2]{%
  \begin{tcolorbox}[colback=purple!4,colframe=purple!45!black,
                    title={Video: \texttt{#1} (wide)}]
  \emph{Caption:} #2
  \end{tcolorbox}}

% the figures inside share one row on the page; on paper they simply follow
% one another
\newenvironment{figrow}{}{}

% two pictures side by side at the text's width, under one caption:
% #1 and #2 file stems, #3 the caption
\newcommand{\figpair}[3]{%
  \begin{tcolorbox}[colback=blue!4,colframe=blue!40!black,
                    title={Figure pair: \texttt{#1} and \texttt{#2}}]
  \emph{Caption:} #3
  \end{tcolorbox}}

% the movie gallery: the environment is the grid, each \movie a poster in it.
% #1 group/stem under movies/ in the data repository, #2 title, #3 caption
\newenvironment{gallery}{}{}
\newcommand{\movie}[3]{%
  \begin{tcolorbox}[colback=purple!4,colframe=purple!45!black,
                    title={Movie: \texttt{#1}}]
  \textbf{#2}: #3
  \end{tcolorbox}}

% a large centred quote. \pullquote{one} on its own, or \pullquote[after]{before}
% for the two-part form with an arrow between them
\makeatletter
\newcommand{\pullquote}{\@ifstar\pq@bare\pq@quoted}
\newcommand{\pq@quoted}[2][]{%
  \begin{center}\itshape\large
  ``#2''\ifx\relax#1\relax\else\ $\Rightarrow$\ ``#1''\fi
  \end{center}}
\newcommand{\pq@bare}[2][]{%
  \begin{center}\itshape\large
  #2\ifx\relax#1\relax\else\ $\Rightarrow$\ #1\fi
  \end{center}}
\makeatother

% the paper this post accompanies. Not printed on the page: it is the source
% of the page's search metadata (the description a search engine shows, and
% the JSON-LD naming the paper, its authors and its abstract)
\newcommand{\papertitle}[1]{\def\thepapertitle{#1}}
\newcommand{\paperauthors}[1]{\def\thepaperauthors{#1}}
\newcommand{\paperaffil}[1]{\def\thepaperaffil{#1}}
\newcommand{\paperabstract}[1]{\def\thepaperabstract{#1}}

\newcommand{\papercard}{%
  \begin{tcolorbox}[colback=cyan!4,colframe=cyan!45!black,
                    title={Page metadata: the paper (not printed on the page)}]
  \textbf{Title:} \thepapertitle\par\smallskip
  \textbf{Authors:} \thepaperauthors
  \begin{enumerate}\itemsep0pt \thepaperaffil \end{enumerate}
  \textbf{Abstract:} \thepaperabstract
  \end{tcolorbox}}

% the blog.html listing: the page's own heading, and the card this post
% gets in the list. Everything here is text on that page, not on the post.
\newcommand{\indexpretitle}[1]{\def\theindexpretitle{#1}}
\newcommand{\indextitle}[1]{\def\theindextitle{#1}}
\newcommand{\postexcerpt}[1]{\def\thepostexcerpt{#1}}
\newcommand{\postthumbalt}[1]{\def\thepostthumbalt{#1}}
\newcommand{\postmore}[1]{\def\thepostmore{#1}}

\newcommand{\indexcard}{%
  \begin{tcolorbox}[colback=green!4,colframe=green!45!black,
                    title={Listing page: \texttt{blog.html}}]
  \textbf{Page heading:} \theindexpretitle{} / \theindextitle

  \medskip
  \textbf{Card}\par
  \textit{Date:} \thepostdate\par
  \textit{Title:} \theposttitle\par
  \textit{Excerpt:} \thepostexcerpt\par
  \textit{Tags:} \theposttags\par
  \textit{Link:} \thepostmore\par
  \textit{Thumbnail alt text:} \thepostthumbalt
  \end{tcolorbox}}

% the downloads picker at the foot of the post
\newcommand{\datalist}[1]{%
  \begin{tcolorbox}[colback=gray!6,colframe=gray!50,
                    title={Generated: downloads picker}]
  #1
  \end{tcolorbox}}

\posttitle{\astrolabe: A 3D Particle-Level Surrogate \\ for Cosmology Simulation \\ with Physical Flow Matching}
\postdate{September 2026}
\posttags{Cosmology; Flow Matching; Generative Models; Simulation Emulation}
\postbyline{by Ibrahim Elsharkawy, last edited September 2026}

\papertitle{\astrolabe: A 3D Particle-Level Surrogate for Cosmology Simulation
with Physical Flow Matching}
\paperauthors{Ibrahim Elsharkawy$^{1,2}$, Yonatan Kahn$^{1}$,
Neal Weiner$^{3}$, David Curtin$^{4}$}
\paperaffil{%
  \item Department of Physics, University of Toronto and Vector Institute,
        Toronto, ON, M5S 1A7, Canada
  \item NERSC, Lawrence Berkeley National Laboratory, Berkeley, California, USA
  \item Center for Cosmology and Particle Physics, Department of Physics,
        New York University, New York, NY 10003, USA
  \item Department of Physics, University of Toronto, Toronto, ON, M5S 1A7,
        Canada}
\paperabstract{}

% the page is unlisted: nothing links to it and it asks not to be indexed,
% so there is no card for it on blog.html yet. These are here for when it
% is published.
\indexpretitle{Blog}
\indextitle{Research Thoughts}
\postexcerpt{}
\postmore{Read post $\rightarrow$}
\postthumbalt{}

\begin{document}
\maketitle
\noindent\textbf{Tags:} \theposttags

% ---------------------------------------------------------------- opening --

% the byline sits under the opening figure, not in the page header
\noindent\textit{\thepostbyline}\par\medskip


\interactive{astro3d}{A 3D interactive figure to inspect N-body Simulation vs \astrolabe. The N-body simulation is on the left of the draggable bar, and \astrolabe is on the right. 
Play runs a pre-designed camera path while $z$ evolves. You can also orbit, zoom and pan yourself at any redshift, and drag the bar across to compare.}

The microphysics of dark matter is thought to be imprinted on cosmological structure across an enormous dynamic range. The most reliable forward model of dark
sector physics to nonlinear cosmological observables is the $N$-body or hydrodynamic simulation. At a cost of $\mathcal{O}(10^6)$ CPU-hours per full-resolution run, systematic scans over dark-matter theory space are computationally intractable with $N$-body simulations. Machine learning emulators provide a possible solution, but are currently limited by either low resolution or only the ability to emulate specific observables. We introduce \astrolabe{}, a flow-matching-based generative emulator designed to be free of both constraints. \astrolabe{} acts directly on the $6N$ phase-space coordinates of individual particles. With no voxel grid in the output, there is no voxel-resolution ceiling, and in principle any observable can be measured from inference. \astrolabe{} is made computationally tractable by employing a field-particle hybrid transformer in which every particle cross-attends to multiscale field tokens, so compute and memory are linear in particle count $N$. We further introduce so-called physical flow matching to improve inductive bias. The flow transports approximate COmoving Lagrangian Acceleration (COLA) \cite{Tassev:2013pn} based positions rather than noise, and momentum conservation is exactly enforced at all points along flow time.

\figplace[medium]{fig-flowtime-decoupled}{Illustration of physical flow
matching, with flow time decoupled from redshift. \textbf{Bottom:} Simulation
initial conditions. \textbf{Middle row:} COLA fields, evolved under cosmic time
from the initial conditions. \textbf{Top row:} \astrolabe{} density field,
evolved under flow time from the COLA distributions in the middle row.}

In the cases we test, \astrolabe{} produces realistic fields with accurate power spectra, halo mass functions, and velocity profiles for the CDM fields it was trained on. We also observe out-of-distribution generalization, with accurate emulation of dark-acoustic oscillation initial conditions. We further find \astrolabe{} capable of being fine-tuned to simulate large particle numbers, and we demonstrate this by emulating a TNG-50-dark simulation at a fraction of the cost. This method represents a first step toward a general and powerful cosmology simulation surrogate, which may eventually permit the inclusion of baryonic physics and richer dark sector interactions.

\figpair{cdm-2m-10mpc-zoom-z0}{mass-profile-1mpc-z0}{\textbf{Left:} Example
CDM field for $128^3\approx2\,\rm M$ particles in a 10 Mpc/$h$ box. $\Sigma$ is
a surface density in units of particles per pixel. \textbf{Top left:} A
\textsc{GIZMO} \cite{Hopkins:2014qka} $N$-body simulation, which serves as the
truth-level phase space density for this example. \textbf{Top right:} Coarse
approximation to the truth-level density given by COLA with a mesh of $256^3$
and 40 timesteps \cite{Tassev:2013pn}. \textbf{Bottom left:} \astrolabe{}-Small
evaluated with ten integration steps and no subsampling. \textbf{Bottom right:}
\astrolabe{}-Small Fast Inference evaluated with 2 integration steps and $c=8$
subsampling. \textbf{Right:} Consistency test for halo density profiles in the
1 Mpc/$h$ box. Halo mass profile averaged over all halos with more than 5000
particles, with concentration $c$ from an NFW fit and inner slope $\gamma$
displayed.}

\section{Gallery}

Movies of N-body Simulation vs \astrolabe for various test cases.  

\subsection{2M CDM}

\begin{gallery}
\movie{2M-CDM/cinematic-2M-10Mpc-CDM-regen23-truth-vs-stage2}{CDM, 10 Mpc/h}{N-body simulation vs Astrolabe-Small}
\movie{2M-CDM/cinematic-2M-1Mpc-CDM-regen23-truth-vs-stage2}{CDM, 1 Mpc/h}{N-body simulation vs Astrolabe-Small}
\end{gallery}

\subsection{2M DAO}

\begin{gallery}
\movie{2M-DAO/cinematic-2M-10Mpc-DAOweak-run1-truth-vs-stage2}{DAO weak, 10 Mpc/h}{$k_{\rm peak} = 31\,h$/Mpc; N-body simulation vs Astrolabe-Small}
\movie{2M-DAO/cinematic-2M-10Mpc-DAOstrong-run1-truth-vs-stage2}{DAO strong, 10 Mpc/h}{$k_{\rm peak} = 18\,h$/Mpc; N-body simulation vs Astrolabe-Small}
\movie{2M-DAO/cinematic-2M-1Mpc-DAOweak-run1-truth-vs-stage2}{DAO weak, 1 Mpc/h}{$k_{\rm peak} = 95\,h$/Mpc; N-body simulation vs Astrolabe-Small}
\movie{2M-DAO/cinematic-2M-1Mpc-DAOstrong-run1-truth-vs-stage2}{DAO strong, 1 Mpc/h}{$k_{\rm peak} = 32\,h$/Mpc; N-body simulation vs Astrolabe-Small}
\end{gallery}

\subsection{TNG-50/100-Dark-3 \cite{Nelson:2018uso}}

\begin{gallery}
\movie{TNG/cinematic-TNG50-3-Dark-z0-truth-vs-stage2}{TNG50-3-Dark, $z = 0$}{N-body simulation vs Astrolabe-Small, fine-tuned on TNG50}
\movie{TNG/cinematic-TNG100-3-Dark-z0-truth-vs-stage2}{TNG100-3-Dark, $z = 0$}{N-body simulation vs Astrolabe-Small, fine-tuned on TNG50}
\end{gallery}

% Body goes here. Sections become <h2> on the page, subsections <h3>,
% paragraphs <p>. See the mapping rules at the top of this file.


% ------------------------------------------- the paper, as page metadata --
\papercard

% ------------------------------------------------------------ references --
\begin{thebibliography}{9}

\bibitem{Tassev:2013pn}
S.~Tassev, M.~Zaldarriaga and D.~Eisenstein, \emph{Solving Large Scale
Structure in Ten Easy Steps with COLA}, JCAP \textbf{06} (2013) 036,
\href{https://arxiv.org/abs/1301.0322}{arXiv:1301.0322}.

\bibitem{Hopkins:2014qka}
P.~F.~Hopkins, \emph{A new class of accurate, mesh-free hydrodynamic simulation
methods}, Mon. Not. Roy. Astron. Soc. \textbf{450} (2015) 53--110,
\href{https://arxiv.org/abs/1409.7395}{arXiv:1409.7395}.

\bibitem{Nelson:2018uso}
D.~Nelson \emph{et al.}, \emph{The IllustrisTNG simulations: public data
release}, Comput. Astrophys. Cosmol. \textbf{6} (2019) 2,
\href{https://arxiv.org/abs/1812.05609}{arXiv:1812.05609}.

% \bibitem{key}
% A.~Author, \emph{Title}, Journal \textbf{vol} (year) pages,
% \href{https://arxiv.org/abs/0000.00000}{arXiv:0000.00000}.

\end{thebibliography}

% the card this post would get on blog.html, once it is listed there
\indexcard

\end{document}
